\documentclass[11pt]{article}
\usepackage[margin=0.78in]{geometry}
\usepackage{amsmath,amssymb,booktabs,tabularx,array,xcolor,hyperref}
\hypersetup{colorlinks=true,urlcolor=blue,linkcolor=blue}
\setlength{\parindent}{0pt}
\setlength{\parskip}{5pt}
\newcolumntype{Y}{>{\raggedright\arraybackslash}X}
\newcommand{\unit}[1]{\,\mathrm{#1}}

\title{S07 reachable near-drop preparation analysis}
\author{Documentary and analytical review artifact}
\date{29 September 2026}

\begin{document}
\maketitle

\section{Question and claim boundary}

The accepted task is a $2{,}400\unit{m}$ three-lane approach followed by a
three-to-two lane reduction.  Common desired-speed advice applies on
$[800,2000]\unit{m}$; the hidden-lane contrast is required in the final
$400\unit{m}$, $R=[2000,2400]\unit{m}$, immediately before the narrowing.
The preparation starts from an empty road and the strong lane-blind channel
retains complete spatial count, density, mean-speed, cumulative-count, source-
waiting, and advice histories through decision time.

This analysis asks when a nontrivial contrast can reach $R$, how a small
outcome-blind bank can avoid obvious insertion confounding, and why exact
aggregate-history equality still cannot be certified without a native reference
run.  It does not supply such a run, approve revised settings, or establish that
lane information changes useful advice.

\section{Arrival bounds}

Let the default base insertion put a $5\unit{m}$ vehicle's front at
$x_0=5\unit{m}$.  Define
\[
x_z=800\unit{m},\qquad x_R=2000\unit{m},\qquad
v_{100}=\frac{100}{3.6}=27.7778\unit{m/s},\qquad
v_{95}=\frac{95}{3.6}=26.3889\unit{m/s},
\]
where $95\unit{km/h}$ is the nominal free speed induced by baseline advice
$90\unit{km/h}$ and response gain $0.5$.  The declared acceleration is
$a=2.6\unit{m/s^2}$.

Two strict kinematic screens ignore all safety, interaction, insertion, and
advice delays.  Instantaneous travel at $v_{100}$ gives
\[
\tau_{\mathrm{inst}}=\frac{x_R-x_0}{v_{100}}=71.8200\unit{s}.
\]
Starting from rest, accelerating to $v_{100}$, and then maintaining it gives
\[
\tau_{\mathrm{acc}}=\frac{v_{100}}{a}+
\frac{x_R-x_0-v_{100}^2/(2a)}{v_{100}}
=77.1619\unit{s}.
\]
These reproduce S06's robust exclusion of the $40\unit{s}$ and $60\unit{s}$
decisions.

For planning rather than proof, two nominal free-flow travel times respect the
$95\unit{km/h}$ advice-zone speed.  Insertion at desired freeway speed gives
\[
\tau_{\mathrm{des}}=
\frac{x_z-x_0}{v_{100}}+\frac{x_R-x_z}{v_{95}}
=74.0937\unit{s}.
\]
Insertion from rest, acceleration to $v_{100}$ before $x_z$, and nominal
$v_{95}$ travel through the advice zone gives
\[
\tau_0=\frac{v_{100}}{a}+
\frac{x_z-x_0-v_{100}^2/(2a)}{v_{100}}
+\frac{x_R-x_z}{v_{95}}=79.4356\unit{s}.
\]
The latter is not a guaranteed upper bound: following, insertion delay, lane
changing, and downstream interaction can make travel later.

After $x_R$, advice no longer lowers desired speed.  Nominal acceleration from
$v_{95}$ to $v_{100}$ gives a region residence estimate
\[
\delta_R=\frac{v_{100}-v_{95}}{a}+
\frac{400-\tfrac12(v_{95}+v_{100})(v_{100}-v_{95})/a}{v_{100}}
=14.4134\unit{s}.
\]

\section{Demand transition and nominal region population}

For intended departure times $d_i$, a free-flow population screen is
\[
N_R(T;\tau,\delta_R)=
\sum_i \mathbf{1}\{T-\tau-\delta_R<d_i\leq T-\tau\}.
\]
The fixed demand schedule has $4\unit{s}$ headway from $0$ to $120\unit{s}$
($900\unit{veh/h}$), then $8/9\unit{s}$ headway from $120$ to $480\unit{s}$
($4050\unit{veh/h}$).  Thus the peak cohort first reaches $R$ nominally at
$194.0937\unit{s}$ under desired-speed insertion and $199.4356\unit{s}$ under
zero-speed insertion.  A full region-residence interval of peak arrivals begins
near $208.5070\unit{s}$ or $213.8489\unit{s}$, respectively.

\begin{table}[h]
\centering
\caption{Nominal vehicles in $R$; these are arithmetic screens, not reference outcomes.}
\begin{tabular}{@{}rcc@{}}
\toprule
Decision $T$ (s) & desired-speed insertion & zero-speed insertion \\
\midrule
80  & 2  & 1 \\
120 & 4  & 4 \\
200 & 9  & 4 \\
210 & 16 & 12 \\
220 & 17 & 17 \\
240 & 16 & 16 \\
260 & 17 & 17 \\
\bottomrule
\end{tabular}
\end{table}

The old $80\unit{s}$ snapshot is sparse even under the faster insertion
convention.  The grid $T\in\{220,240,260\}\unit{s}$ lies beyond the nominal
peak-fill time under both conventions and preserves three stages for a later
prefix screen.  Actual counts remain audit outputs, not acceptance assumptions.

\section{Why the old quota bank is poorly conditioned}

At peak demand, assigning $k$ of every six intended vehicles to the ending lane
creates the following nominal lane rates; the two continuing-lane rates may swap
with phase.

\begin{table}[h]
\centering
\begin{tabular}{@{}rccc@{}}
\toprule
Ending quota $k/6$ & ending lane & continuing lane A & continuing lane B \\
\midrule
1/6 & 675  & 2025 & 1350 \\
2/6 & 1350 & 1350 & 1350 \\
3/6 & 2025 & 1350 & 675  \\
4/6 & 2700 & 675  & 675  \\
\bottomrule
\end{tabular}
\end{table}

Units are vehicles per hour.  SUMO's official capacity examples show that
insertion throughput depends strongly on departure speed, step length, and
placement; zero departure speed is specifically the weakest example.  The table
above is therefore not compared to a universal capacity threshold.  It does show
that the unbalanced quotas impose very different lane-specific insertion burdens.
Because source waiting and actual cumulative entries are part of the strong
history, preparation-dependent insertion delay can reveal the preparation before
the lane-drop mechanism is tested.

The cleanest small-bank structure uses one common balanced background and changes
only one six-vehicle contrast block.  All six permutations below give two
departures to each lane and separate same-lane departures by three intended
headways:
\[
\begin{array}{lll}
P_0=(0,1,2,0,1,2), & P_1=(0,2,1,0,2,1),\\
P_2=(1,0,2,1,0,2), & P_3=(2,0,1,2,0,1),\\
P_4=(1,2,0,1,2,0), & P_5=(2,1,0,2,1,0).
\end{array}
\]
Lane $0$ is the ending lane.  Outside the contrast block every preparation uses
$P_0$; from decision time onward every preparation also uses the same frozen
pattern.  This localizes the hidden input and holds the block's per-lane totals
fixed, but it does not guarantee equal actual insertion or trajectories.

Let peak index $j=0$ depart at $120\unit{s}$ with headway $8/9\unit{s}$.
For the three proposed decisions, the following complete six-vehicle blocks fall
inside $R$ at decision under both nominal travel conventions:

\begin{table}[h]
\centering
\begin{tabular}{@{}rcc@{}}
\toprule
$T$ (s) & peak indices $j$ & intended departure interval (s) \\
\midrule
220 & 16--21 & 134.222--138.667 \\
240 & 38--43 & 153.778--158.222 \\
260 & 61--66 & 174.222--178.667 \\
\bottomrule
\end{tabular}
\end{table}

This yields $3\times6=18$ prefix preparations, down from the old
$3\times16=48$, without selecting by action outcome.

\section{Anticipation and exact observability}

The native lane-change model prevents an analytical equality certificate.
Erdmann's primary LC2013 description states that strategic lane changing may
occur well before a dead lane and that urgency depends on remaining distance,
current/historical speed, lane offset, and target-lane occupation.  Blocked
changes may alter the ego and neighboring vehicles' speeds.  Those changes feed
the exact per-step counts and mean-speed histories retained by the comparator.

Current official SUMO documentation lists a default
\texttt{lcStrategicLookahead} of $3000\unit{m}$ and a left-change lookahead
factor of $2$.  This rolling documentation is not a claim about an uninspected
future execution version, but its $3000\unit{m}$ scale exceeds the entire
$2400\unit{m}$ approach.  Consequently, it is not defensible to assume a long
upstream interval in which ending-lane assignments are dynamically irrelevant.
Nor may S07 shorten lookahead or disable strategic changes merely to manufacture
an exact collision.

The localized balanced block therefore improves the preparation audit without
proving success.  A later native prefix run must first require exact equality of
the declared arrays and metadata.  Any mismatch stops the action test.  A failed
bank would show that this contract did not produce a collision; it would not show
that aggregate histories are universally sufficient.

\section{Preparation-family implications}

\textbf{Extended road-start warm-up.}  Keep the accepted demand clock.  Set
\texttt{departPos} to \texttt{base} and \texttt{departSpeed} to
\texttt{desired}; retain all
insertion safety and delays, use $T=220,240,260\unit{s}$, and use the localized
balanced block.  This is the smallest repair and is the provisional
recommendation.  Its exact-match probability remains unknown.

\textbf{Shift peak demand earlier.}  Starting the balanced peak at time zero can
populate $R$ near $90\unit{s}$, reducing warm-up length.  It changes the demand
regime and removes the original low-demand lead-in without reducing strategic
anticipation or exact-observability risk.  It is therefore a broader amendment
with no established scientific advantage.

\textbf{Closer or nonempty initialization.}  Numeric mid-link departure positions
are supported by the reference interface and could create a lane-blind collision
immediately.  In this task they introduce an artificial near-drop source or a
hidden initial state, changing the known-empty/full-history contract and shifting
the physical-reachability burden outside the model.  This is not an admissible
Candidate A repair unless DG-G1/DG-G2 are explicitly reopened.

\section{Source boundary}

Primary/native basis: J. Erdmann, ``SUMO's Lane-Changing Model,'' 2015,
\href{https://elib.dlr.de/102254/1/Springer-SUMOs_Lane_changing_model.pdf}{DLR manuscript};
SUMO official
\href{https://sumo.dlr.de/docs/Definition_of_Vehicles,_Vehicle_Types,_and_Routes.html}{vehicle/type/route documentation},
\href{https://sumo.dlr.de/docs/Simulation/VehicleInsertion.html}{vehicle insertion documentation}, and
\href{https://sumo.dlr.de/docs/Simulation/RoadCapacity.html}{road-capacity methodology}.
Rolling documentation supports capability and design-risk statements only; no
installed version or empirical scenario result is inferred.

\end{document}
